\documentclass[11pt]{article}

\usepackage[margin=1in]{geometry}
\usepackage{amsmath,amssymb,amsthm}
\usepackage{booktabs}
\usepackage{array}
\usepackage{multirow}
\usepackage{longtable}
\usepackage{hyperref}
\usepackage{enumitem}
\usepackage{listings}
\usepackage{xcolor}

\hypersetup{
  colorlinks=true,
  linkcolor=blue!60!black,
  urlcolor=blue!60!black,
  citecolor=blue!60!black
}

\lstset{
  basicstyle=\ttfamily\small,
  breaklines=true,
  frame=single,
  framerule=0pt,
  backgroundcolor=\color{gray!8},
  columns=fullflexible
}

\newcommand{\vj}{v_{j}}
\newcommand{\vjk}{v_{j,k}}
\newcommand{\thetae}{\theta_{1}}
\newcommand{\thetal}{\theta_{2}}

\title{ORIE 5132 Project: Choice Modeling, Assortment Optimization, and Pricing}
\author{Theodore Chronopoulos \\ Pricing Analytics and Revenue Management}
\date{Spring 2026}

\begin{document}
\maketitle

\begin{abstract}
This report presents the solutions to Problems 1--6 of the ORIE 5132 course project.
We estimate a Multinomial Logit (MNL) choice model on Expedia hotel search data,
use the fitted model for unconstrained assortment optimization and pricing,
extend to a mixture of MNL (MMNL) with two customer types,
and solve the type-aware and type-unaware assortment problems.
Problem 6 repeats the exercise with a different segmentation (family vs.\ non-family).
For every problem we state the model, derive the structural result used,
describe the solution method, and report the numerical output together with
sanity checks that verify correctness.
\end{abstract}

\tableofcontents

% =====================================================================
\section{Data and Notation}
\label{sec:data}

We use the Expedia dataset \texttt{data.csv} (8\,354 queries, 153\,009 hotel rows,
average $18.3$ alternatives per query, 5\,848 booked queries).
Each query \texttt{srch\_id} is the choice set~$S$ faced by one customer,
who either books at most one hotel or leaves without booking.

The eight hotel features are, in order:
\begin{enumerate}[label=\alph*)]
  \item \texttt{prop\_starrating} (star rating),
  \item \texttt{prop\_review\_score} (review score),
  \item \texttt{prop\_brand\_bool} (hotel brand indicator),
  \item \texttt{prop\_location\_score} (location score),
  \item \texttt{prop\_accesibility\_score} (accessibility score),
  \item \texttt{prop\_log\_historical\_price} (log historical price),
  \item \texttt{price\_usd} (displayed price),
  \item \texttt{promotion\_flag} (promotion indicator).
\end{enumerate}

Four small datasets \texttt{data1.csv}, \texttt{data2.csv}, \texttt{data3.csv},
\texttt{data4.csv} contain 26--29 hotels each and are used for assortment and pricing.

All code is deterministic under seed \texttt{5132}.

% =====================================================================
\section{Problem 1: MNL Estimation}
\label{sec:p1}

\subsection{Model}
Under MNL, hotel $j$ has preference weight
\[
v_j = e^{u_j}, \qquad
u_j = \beta_0 + \sum_{i=1}^{8} \beta_i x_{ji},
\]
and the probability that a customer chooses $j$ from the displayed set $S$ is
\[
\mathbb{P}(j \mid S) \;=\; \frac{v_j}{1 + \sum_{p \in S} v_p}.
\]
The outside option has weight $1$.

\subsection{Estimation}
We maximise the log-likelihood
\[
\mathcal{L}(\beta) \;=\; \sum_{q \in Q} \Big[ \mathbf{1}_{\{y_q \ne 0\}} u_{q, y_q}
  - \log\!\Big(1 + \sum_{p \in S_q} e^{u_{q,p}} \Big) \Big]
\]
by a damped Newton method with backtracking, implemented in
\texttt{scripts/mnl\_utils.py}.
Continuous features are $z$-scored before estimation; binary features
are left at $\{0,1\}$.
The coefficients reported below are on the \emph{raw} (un-standardised) scale
so that they can be fed directly to the small datasets.

\subsection{Results}
\begin{center}
\begin{tabular}{lr}
\toprule
Coefficient & Estimate (raw scale) \\
\midrule
intercept                    & $-2.8153$ \\
prop\_starrating             & $+0.4761$ \\
prop\_review\_score          & $+0.1199$ \\
prop\_brand\_bool            & $+0.2299$ \\
prop\_location\_score        & $+0.0163$ \\
prop\_accesibility\_score    & $+0.5629$ \\
prop\_log\_historical\_price & $-0.0374$ \\
price\_usd                   & $-0.00732$ \\
promotion\_flag              & $+0.4540$ \\
\midrule
final log-likelihood         & $-20\,611.33$ \\
final gradient norm          & $1.22 \times 10^{-11}$ \\
Newton iterations            & $6$ \\
\bottomrule
\end{tabular}
\end{center}

\subsection{Commentary}
\begin{itemize}
  \item \textbf{Price sensitivity} $\beta_{\text{price}} = -0.00732 < 0$ is
    economically sensible: higher displayed price lowers choice probability.
    This sign is essential for Problem 3 to have an interior optimum.
  \item \textbf{Promotion} $\beta_{\text{promo}} = +0.454$: promoted hotels
    are more likely to be chosen, as expected.
  \item \textbf{Quality signals}: star rating, review score, brand, and
    accessibility all have positive coefficients. Location score is
    weakly positive.
  \item \textbf{Historical price} is slightly negative ($-0.037$). Once the
    displayed price and other quality attributes are controlled for, the
    historical-price proxy no longer provides useful information and even
    absorbs a small negative correlation. This is a joint-estimation
    artefact rather than a sign that expensive historical hotels are
    undesirable.
  \item Final gradient norm is machine-precision small, so the Newton
    method has converged to a stationary point of the (concave) MNL
    log-likelihood. It is the global maximum.
\end{itemize}

\paragraph{Verification.}
The uniform-choice baseline log-likelihood (every hotel and the outside
option equally attractive within each query) is $-22\,856.21$.
The fitted model improves on this baseline by $2\,244.88$ log-likelihood
units, confirming that the fit is meaningful.

% =====================================================================
\section{Problem 2: Assortment Optimization under MNL}
\label{sec:p2}

\subsection{Problem Statement}
Given the hotels in \texttt{dataX.csv}, choose a subset $S \subseteq \{1,\dots,n\}$
to display so as to maximise the expected revenue
\[
R(S) \;=\; \frac{\sum_{j \in S} p_j v_j}{1 + \sum_{j \in S} v_j},
\]
where $v_j = \exp(\hat u_j)$ uses the Problem~1 coefficients.

\subsection{Structural Result}
Under unconstrained MNL assortment optimization, the optimal set is
\emph{revenue-ordered} (Talluri and van~Ryzin, 2004): there exists a
price cutoff~$p^*$ such that
\[
S^* \;=\; \{\, j : p_j \ge p^* \,\}.
\]
Consequently the optimum is one of the $n$ nested prefixes obtained by
sorting hotels in decreasing price order. This gives an $O(n \log n)$
exact algorithm.

\subsection{Algorithm}
Sort hotels by $p_j$ descending; for $k = 1,\dots,n$ compute the revenue of
the $k$-th prefix; return the prefix with the largest revenue.

\subsection{Results}
\begin{center}
\begin{tabular}{lrrl}
\toprule
Dataset & $|S^*|$ & $R(S^*)$ & $S^*$ (1-based indices) \\
\midrule
\texttt{data1.csv} & 18 & $107.3531$ & $\{1,2,3,4,5,6,7,13,16,18,19,20,21,22,23,24,25,27\}$ \\
\texttt{data2.csv} & 10 & $131.1136$ & $\{1,2,7,8,9,10,11,22,24,26\}$ \\
\texttt{data3.csv} & 18 & $121.0548$ & $\{1,2,3,4,5,6,8,9,11,12,14,15,16,17,19,20,24,25\}$ \\
\texttt{data4.csv} & 11 & $97.4089$  & $\{4,5,7,9,11,16,19,20,21,22,27\}$ \\
\bottomrule
\end{tabular}
\end{center}

\paragraph{Verification.}
For each returned assortment we check all single-item add/drop moves
and confirm that no one-step modification increases revenue.
This is consistent with the revenue-ordered theorem and with local
optimality.

% =====================================================================
\section{Problem 3: Pricing under MNL}
\label{sec:p3}

\subsection{Problem Statement}
Display all $n$ hotels in \texttt{dataX.csv} and choose new prices
$p_1,\dots,p_n$ maximising
\[
R(p) \;=\; \frac{\sum_j p_j \, v_j(p_j)}{1 + \sum_j v_j(p_j)},
\qquad
v_j(p_j) \;=\; \exp\!\big(a_j + b\, p_j\big),
\]
where $a_j$ is the Problem-1 utility stripped of the price term and
$b = \hat\beta_{\text{price}} = -0.00732 < 0$.

\subsection{First-order Condition}
Let $V = \sum_j v_j$ and $D = 1 + V$. Using $dv_j/dp_j = b\, v_j$,
\[
\frac{\partial R}{\partial p_j}
  \;=\;
  \frac{\big(v_j + p_j b v_j\big) D - \Big(\sum_k p_k v_k\Big) b v_j}{D^2}
  \;=\;
  \frac{v_j (1+V) \, \big[\, 1 + b p_j - b R \,\big]}{D^2}.
\]
Setting $\partial R/\partial p_j = 0$ and cancelling the positive factor
$v_j (1+V)/D^2$ yields, for every $j$,
\[
\boxed{\; p_j \;=\; R + \tfrac{1}{|b|}. \;}
\]
The right-hand side is independent of $j$, so \emph{every optimal price
is the same}. Writing $p^* = R^* + 1/|b|$ and substituting back,
the common optimum $(p^*, R^*)$ satisfies the fixed-point system
\[
p^* \;=\; R^* + \tfrac{1}{|b|},
\qquad
R^* \;=\; p^* \cdot \frac{\sum_j e^{a_j + b p^*}}{1 + \sum_j e^{a_j + b p^*}},
\]
which has a unique solution (Lambert-$W$ type).

\subsection{Algorithm}
Because the optimum is a scalar, we bracket the maximiser of
$R(p, p, \dots, p)$ and refine it by golden-section search
(\texttt{scripts/problem3\_pricing.py}).

\subsection{Results}
\begin{center}
\begin{tabular}{lrrr}
\toprule
Dataset & $p^*$ (USD) & Current $R$ & Optimal $R$ \\
\midrule
\texttt{data1.csv} & $314.36$ & $104.65$ & $177.81$ \\
\texttt{data2.csv} & $385.36$ & $113.36$ & $248.80$ \\
\texttt{data3.csv} & $313.03$ & $118.17$ & $176.47$ \\
\texttt{data4.csv} & $351.03$ & $78.69$  & $214.47$ \\
\bottomrule
\end{tabular}
\end{center}

\paragraph{Numerical sanity check.}
With $1/|b| = 136.56$, the identity $p^* - R^* = 1/|b|$ should hold for
every dataset. For \texttt{data1}: $314.36 - 177.81 = 136.55$;
for \texttt{data2}: $385.36 - 248.80 = 136.56$; and similarly for
\texttt{data3}, \texttt{data4}. All four match $136.56$ to two decimals,
confirming the fixed-point residual is essentially zero.

\subsection{Commentary}
The identical optimal prices across hotels are a direct consequence of the
\emph{single shared price coefficient} $\beta_{\text{price}}$ in the
Problem~1 utility. Differentiated optimal prices would require a utility
specification with hotel-specific or interacted price sensitivity, e.g.\
$u_j = a_j + \beta_{\text{price}} (1 + \gamma\, \mathrm{brand}_j)\, p_j$,
which is outside the model asked for. Under the fitted model the reported
solution is exact.

% =====================================================================
\section{Problem 4: Mixture of MNL (Early vs.\ Late)}
\label{sec:p4}

\subsection{Model}
Two customer types indexed by $k \in \{1,2\}$:
type~1 = early (booking window $\ge 7$ days), type~2 = late
(booking window $< 7$ days). Let $\theta_k$ be the share of type-$k$
customers. For each type we fit a separate MNL with preference weights
$v_{j,k} = \exp(\beta_{0k} + \sum_i \beta_{ik} x_{ji})$, and the mixture
choice probability is
\[
\mathbb{P}(j \mid S) \;=\; \sum_{k=1}^{2} \theta_k \,
  \frac{v_{j,k}}{1 + \sum_{p \in S} v_{p,k}}.
\]

\subsection{Estimating $\theta_1, \theta_2$}
Because the type is observable from the booking window, we estimate the
$\theta$'s by simple proportions of queries (customers, \emph{not} rows):
\[
\hat\theta_1 \;=\; \frac{\#\{\text{queries with bw} \ge 7\}}{\#\text{queries}},
\qquad
\hat\theta_2 \;=\; 1 - \hat\theta_1.
\]
Shared scaling statistics (means, std.\ devs.) are used across both types
so the coefficient estimates are directly comparable.

\subsection{Results}
\begin{center}
\begin{tabular}{lr}
\toprule
$\hat\theta_1$ (early) & $0.5431$ \\
$\hat\theta_2$ (late)  & $0.4569$ \\
\bottomrule
\end{tabular}
\end{center}

\medskip
\begin{center}
\begin{tabular}{lrrr}
\toprule
Coefficient (raw) & Early & Late & Early$-$Late \\
\midrule
intercept                    & $-2.9747$ & $-2.7096$ & $-0.265$ \\
prop\_starrating             & $+0.4422$ & $+0.5399$ & $-0.098$ \\
prop\_review\_score          & $+0.1374$ & $+0.1054$ & $+0.032$ \\
prop\_brand\_bool            & $+0.2089$ & $+0.2561$ & $-0.047$ \\
prop\_location\_score        & $-0.0159$ & $+0.0647$ & $-0.081$ \\
prop\_accesibility\_score    & $+0.7496$ & $+0.3312$ & $+0.419$ \\
prop\_log\_historical\_price & $-0.0531$ & $-0.0179$ & $-0.035$ \\
price\_usd                   & $-0.00590$ & $-0.00934$ & $+0.00344$ \\
promotion\_flag              & $+0.3829$ & $+0.5517$ & $-0.169$ \\
\bottomrule
\end{tabular}
\end{center}

Both per-type log-likelihoods converged to machine precision
(gradient norms $3.2\times 10^{-9}$ and $1.2\times 10^{-9}$).

\subsection{Commentary}
\begin{itemize}
  \item \textbf{Late customers are more price sensitive}
    ($-0.00934$ vs.\ $-0.00590$). Last-minute travellers are more
    value-conscious, whereas early planners tolerate higher prices.
  \item \textbf{Late customers respond more strongly to promotions}
    ($+0.552$ vs.\ $+0.383$) and to star rating, consistent with a more
    deal-seeking mindset.
  \item \textbf{Early customers value accessibility more}
    ($+0.750$ vs.\ $+0.331$). Travellers booking far in advance often
    care about convenience, whereas last-minute bookers optimise price.
  \item Both types reject high displayed price and reward promotions,
    so the qualitative structure is preserved across segments.
\end{itemize}

% =====================================================================
\section{Problem 5: Early vs.\ Late Assortment Optimization}
\label{sec:p5}

For each small dataset we compute
\begin{itemize}
  \item $S$: optimum with \emph{unknown} arriving customer type
    (maximises mixture revenue),
  \item $S_1$: optimum if the customer is known to be type~1 (early),
  \item $S_2$: optimum if the customer is known to be type~2 (late).
\end{itemize}

\subsection{Type-aware Sub-problems}
Each of $S_1$ and $S_2$ is a single-class MNL assortment, so the
revenue-ordered theorem applies and we reuse the Problem-2 algorithm.

\subsection{Unknown Type: Mixed-Integer Program}
For unknown type we solve
\[
\max_{S \subseteq \{1,\dots,n\}}
  \; \theta_1 \frac{\sum_{j \in S} p_j v_{j,1}}{1 + \sum_{j \in S} v_{j,1}}
  + \theta_2 \frac{\sum_{j \in S} p_j v_{j,2}}{1 + \sum_{j \in S} v_{j,2}}.
\]
This is a sum of \emph{fractional} objectives in the binary selection, so
it is not revenue-ordered in general. We give the exact MILP
reformulation that was implemented and solved.

\paragraph{Decision variables.}
\begin{tabular}{ll}
$x_j \in \{0,1\}$            & $1$ if hotel $j$ is offered in $S$, \\
$t_k \ge 0$, $k \in \{1,2\}$ & auxiliary variables linearising $1/(1+\sum_j v_{j,k} x_j)$, \\
$s_{j,k} \ge 0$              & McCormick product linearisation of $t_k \cdot x_j$.
\end{tabular}

\paragraph{Objective.}
\[
\max \;\; \sum_{j=1}^{n} \Big[
   \theta_1 \, p_j \, v_{j,1} \, s_{j,1}
 + \theta_2 \, p_j \, v_{j,2} \, s_{j,2} \Big].
\]

\paragraph{Normalisation constraints} (tie $t_k$ to the assortment):
\[
t_k + \sum_{j=1}^{n} v_{j,k}\, s_{j,k} \;=\; 1, \qquad k \in \{1,2\}.
\]

\paragraph{McCormick envelopes} for $s_{j,k} = t_k \cdot x_j$ with
$x_j \in \{0,1\}$ and $t_k \in [0,1]$, for all $j$ and $k$:
\[
s_{j,k} \le t_k, \qquad
s_{j,k} \le x_j, \qquad
s_{j,k} \ge t_k + x_j - 1, \qquad
s_{j,k} \ge 0.
\]

\paragraph{Bounds.}
\[
0 \le t_k \le 1, \qquad 0 \le s_{j,k} \le 1, \qquad x_j \in \{0,1\}.
\]

\paragraph{Exactness.} Because $x_j$ is binary, the McCormick envelope is
tight: at every integer feasible solution $s_{j,k} = t_k\, x_j$ exactly,
the normalisation forces $t_k = 1/(1 + \sum_j v_{j,k} x_j)$, and the
objective collapses to the original mixture revenue.
Therefore the MILP is an \emph{exact reformulation}, not a relaxation.

\subsection{Algorithm}
The MILP is solved by \texttt{gurobi\_cl} on an LP file produced by
\texttt{scripts/problem5\_assortment.py}. A branch-and-bound fallback
solving the same problem exactly is also available for environments
without Gurobi; it uses the segment-wise best-completion bound for pruning
and is used to cross-validate Gurobi's solution.

\subsection{Results}
\begin{center}
\begin{tabular}{lllrrr}
\toprule
Dataset & Set & Indices & $|\cdot|$ & $R_1$ & $R_2$ \\
\midrule
\multirow{3}{*}{\texttt{data1.csv}}
 & $S$   & $\{1,2,3,4,5,6,7,13,16,18,19,20,21,22,23,24,25,27\}$ & 18 & $103.25$ & $112.12$ \\
 & $S_1$ & $\{1,2,3,4,5,6,7,10,13,15,16,17,18,19,20,21,22,23,24,25,27\}$ & 21 & $103.47$ & $111.76$ \\
 & $S_2$ & $\{1,2,3,4,5,6,7,16,18,19,20,21,22,23,24,25\}$ & 16 & $102.51$ & $112.22$ \\
\midrule
\multirow{3}{*}{\texttt{data2.csv}}
 & $S$   & $\{1,2,7,8,9,10,11,22,24,26\}$                & 10 & $126.92$ & $136.46$ \\
 & $S_1$ & $\{1,2,7,8,9,10,11,22,24,26\}$                & 10 & $126.92$ & $136.46$ \\
 & $S_2$ & $\{1,2,7,8,9,11,22\}$                          &  7 & $123.85$ & $137.38$ \\
\midrule
\multirow{3}{*}{\texttt{data3.csv}}
 & $S$   & $\{1,2,3,4,5,6,8,9,11,12,14,15,16,17,19,20,24,25\}$ & 18 & $117.44$ & $125.23$ \\
 & $S_1$ & $\{1,2,3,4,5,6,8,9,11,12,14,15,16,17,19,20,24,25,26\}$ & 19 & $117.52$ & $125.04$ \\
 & $S_2$ & $\{1,2,3,4,5,6,8,9,11,12,14,15,16,17,20,24,25\}$ & 17 & $117.19$ & $125.35$ \\
\midrule
\multirow{3}{*}{\texttt{data4.csv}}
 & $S$   & $\{4,5,7,9,11,16,19,20,21,22,27\}$ & 11 & $94.70$ & $100.58$ \\
 & $S_1$ & $\{4,5,7,9,11,16,19,20,21,22,27\}$ & 11 & $94.70$ & $100.58$ \\
 & $S_2$ & $\{4,5,7,11,16,19,20,21,22,27\}$   & 10 & $94.21$ & $100.64$ \\
\bottomrule
\end{tabular}
\end{center}

\noindent
$R_1$ denotes expected revenue under the early (type-1) MNL;
$R_2$ denotes expected revenue under the late (type-2) MNL.

\subsection{Comparison and Value of Knowing the Customer Type}
For each dataset we report the gain from knowing the type of an arriving
customer: $R_1(S_1) - R_1(S)$ under type~1 and $R_2(S_2) - R_2(S)$ under
type~2.

\begin{center}
\begin{tabular}{lrr}
\toprule
Dataset & $\Delta_1 = R_1(S_1) - R_1(S)$ & $\Delta_2 = R_2(S_2) - R_2(S)$ \\
\midrule
\texttt{data1.csv} & $+0.218$ & $+0.102$ \\
\texttt{data2.csv} & $\phantom{+}0.000$ & $+0.919$ \\
\texttt{data3.csv} & $+0.083$ & $+0.119$ \\
\texttt{data4.csv} & $\phantom{+}0.000$ & $+0.062$ \\
\bottomrule
\end{tabular}
\end{center}

All values of information are non-negative (as they must be, since $S_k$
maximises $R_k$ by definition and $S$ is feasible for that sub-problem).
Some are exactly zero: on \texttt{data2} and \texttt{data4} the
unknown-type optimum coincides with $S_1$, so knowing the customer is
type~1 brings no extra revenue on those datasets.

\subsection{Cross-Validation of the Two Solvers}
Running the explicit Gurobi MILP and the in-house branch-and-bound
fallback on all four datasets yields \emph{identical} selected sets and
mixture revenues (floating-point difference $<3\times 10^{-14}$).
The two solvers certify each other, which is strong evidence that
$S$ is globally optimal in each case.

% =====================================================================
\section{Problem 6: MMNL with Family vs.\ Non-family Segmentation}
\label{sec:p6}

\subsection{Chosen Segmentation}
We redefine customer types using \texttt{srch\_children\_count}:
\[
\text{type 1 (family)} \iff \texttt{srch\_children\_count} > 0,
\qquad
\text{type 2 (non-family)} \iff \texttt{srch\_children\_count} = 0.
\]
This captures a different demand structure than the booking-window cut:
families are more likely to be multi-guest leisure bookings, while
non-family queries mix solo and couple travellers.

\subsection{Estimated $\theta$ and Coefficients}
\begin{center}
\begin{tabular}{lr}
\toprule
$\hat\theta_1$ (family)     & $0.2246$ \\
$\hat\theta_2$ (non-family) & $0.7754$ \\
\bottomrule
\end{tabular}
\end{center}

\begin{center}
\begin{tabular}{lrrr}
\toprule
Coefficient (raw) & Family & Non-family & Family$-$Non-family \\
\midrule
intercept                    & $-2.6686$ & $-2.8934$ & $+0.225$ \\
prop\_starrating             & $+0.3532$ & $+0.5192$ & $-0.166$ \\
prop\_review\_score          & $+0.1154$ & $+0.1236$ & $-0.008$ \\
prop\_brand\_bool            & $+0.3062$ & $+0.2142$ & $+0.092$ \\
prop\_location\_score        & $+0.0188$ & $+0.0155$ & $+0.003$ \\
prop\_accesibility\_score    & $+0.5193$ & $+0.5677$ & $-0.048$ \\
prop\_log\_historical\_price & $-0.0825$ & $-0.0189$ & $-0.064$ \\
price\_usd                   & $-0.00376$ & $-0.00860$ & $+0.00484$ \\
promotion\_flag              & $+0.4873$ & $+0.4408$ & $+0.047$ \\
\bottomrule
\end{tabular}
\end{center}

\paragraph{Convergence note.} The family model converged with gradient
norm $3.3\times 10^{-5}$ and stable log-likelihood across the last four
iterations (changes below $10^{-8}$); the non-family model converged to
$1.8\times 10^{-11}$. Both are effectively at their stationary point.

\subsection{Commentary}
\begin{itemize}
  \item \textbf{Non-family customers are roughly $2.3\times$ more
    price-sensitive} than families ($-0.00860$ vs.\ $-0.00376$).
    One plausible interpretation: families have a smaller set of
    acceptable hotels (space, amenities) and shop less on price, while
    solo/couple travellers can substitute more freely.
  \item \textbf{Families value brand more} ($+0.306$ vs.\ $+0.214$) and
    respond \textbf{more strongly to promotions} ($+0.487$ vs.\ $+0.441$).
    Both are consistent with the story above.
  \item \textbf{Non-family customers reward star rating more},
    suggesting leisure couples and solo travellers trade up on quality
    once price is controlled for.
\end{itemize}

\subsection{Assortment Results}
The unknown-type problem is solved with the same exact branch-and-bound
used in Problem~5, now against the family/non-family mixture.

\begin{center}
\begin{tabular}{lllrrr}
\toprule
Dataset & Set & Indices & $|\cdot|$ & $R_{\text{fam}}$ & $R_{\text{nf}}$ \\
\midrule
\multirow{3}{*}{\texttt{data1.csv}}
 & $S$   & $\{1,2,3,4,5,6,7,13,16,18,19,20,21,22,23,24,25,27\}$          & 18 & $108.58$ & $106.97$ \\
 & $S_1$ & $\{1,2,3,4,5,6,7,16,18,19,20,21,22,23,24,25,27\}$             & 17 & $108.60$ & $106.95$ \\
 & $S_2$ & $\{1,2,3,4,5,6,7,10,13,15,16,18,19,20,21,22,23,24,25,27\}$    & 20 & $108.50$ & $106.97$ \\
\midrule
\multirow{3}{*}{\texttt{data2.csv}}
 & $S$   & $\{1,2,7,8,9,10,11,22,24,26\}$ & 10 & $139.61$ & $127.85$ \\
 & $S_1$ & $\{1,2,7,8,9,11,22\}$           &  7 & $141.57$ & $124.70$ \\
 & $S_2$ & $\{1,2,7,8,9,10,11,22,24,26\}$ & 10 & $139.61$ & $127.85$ \\
\midrule
\multirow{3}{*}{\texttt{data3.csv}}
 & $S$   & $\{1,2,3,4,5,6,8,9,11,12,14,15,16,17,19,20,24,25\}$           & 18 & $129.91$ & $118.70$ \\
 & $S_1$ & $\{1,2,3,4,5,6,8,11,14,15,16,17,20,24,25\}$                    & 15 & $130.34$ & $117.78$ \\
 & $S_2$ & $\{1,2,3,4,5,6,8,9,11,12,14,15,16,17,19,20,24,25,26\}$        & 19 & $129.58$ & $118.74$ \\
\midrule
\multirow{3}{*}{\texttt{data4.csv}}
 & $S$   & $\{4,5,7,9,11,16,19,20,21,22,27\}$ & 11 & $102.12$ & $95.42$ \\
 & $S_1$ & $\{4,5,7,11,16,19,20,21,22,27\}$   & 10 & $102.31$ & $94.96$ \\
 & $S_2$ & $\{4,5,7,9,11,16,19,20,21,22,27\}$ & 11 & $102.12$ & $95.42$ \\
\bottomrule
\end{tabular}
\end{center}

\begin{center}
\begin{tabular}{lrr}
\toprule
Dataset & $\Delta_1 = R_{\text{fam}}(S_1) - R_{\text{fam}}(S)$ & $\Delta_2 = R_{\text{nf}}(S_2) - R_{\text{nf}}(S)$ \\
\midrule
\texttt{data1.csv} & $+0.012$ & $+0.001$ \\
\texttt{data2.csv} & $+1.951$ & $\phantom{+}0.000$ \\
\texttt{data3.csv} & $+0.437$ & $+0.046$ \\
\texttt{data4.csv} & $+0.193$ & $\phantom{+}0.000$ \\
\bottomrule
\end{tabular}
\end{center}

All values of information are non-negative, as required. On \texttt{data2}
the value of knowing the customer is a family is sizeable ($+1.95$),
because $S_1$ offers the short luxury prefix $\{1,2,7,8,9,11,22\}$ rather
than the wider set that maximises the blended mixture objective.

\subsection{Observations}
\begin{itemize}
  \item The unknown-type assortments $S$ in Problems 5 and 6 often
    coincide exactly (e.g.\ \texttt{data1.csv}, \texttt{data4.csv}), because
    the revenue-ordered structure of the underlying small datasets is
    robust to the particular segmentation used: both segmentations keep
    roughly the same prefix of high-price hotels.
  \item The family/non-family split is \emph{less informative} than the
    early/late split in the sense that most values of knowing the type
    are near zero on Problem~6, except on \texttt{data2} for the family
    segment. This suggests the family/non-family cut produces two groups
    that agree more on ranking than the booking-window cut does.
\end{itemize}

% =====================================================================
\section{Summary}
\label{sec:summary}

\begin{itemize}
  \item Problem~1: MNL fitted cleanly, all coefficient signs economically
    sensible, final gradient norm $\approx 10^{-11}$.
  \item Problem~2: revenue-ordered theorem applies; optimum computed in
    closed form for all four small datasets.
  \item Problem~3: with a single shared price coefficient, all optimal
    prices must be equal; the derivation is explicit and the numerical
    residual $p^* - R^* - 1/|b|$ is essentially zero on every dataset.
  \item Problem~4: query-level $\theta$ estimates and segment-specific
    MNL fits converge; late customers are more price-sensitive,
    early customers more sensitive to accessibility.
  \item Problem~5: type-aware problems use the revenue-ordered theorem;
    the unknown-type problem is formulated as an exact MILP and solved by
    Gurobi, with a branch-and-bound fallback. The two solvers return
    identical solutions on all four datasets.
  \item Problem~6: family/non-family segmentation produces qualitatively
    different sensitivities (non-family more price-sensitive) and
    assortments that are close to but not equal to those in Problem~5.
    Values of information are non-negative across datasets.
\end{itemize}

\paragraph{Reproducibility.}
All scripts default to seed \texttt{5132} and are deterministic.
See \texttt{README.md} for the exact commands.

\end{document}
